\section{AHSME 1953}
        \begin{enumerate}
        	\item (\textit{AHSME 1953}) A boy buys oranges at $3$ for $10$ cents. He will sell them at $5$ for $20$ cents. In order to make a profit of $ $1.00$, he must sell: 


 $\textbf{(A)}\ 67 \text{ oranges} \qquad \textbf{(B)}\ 150 \text{ oranges} \qquad \textbf{(C)}\ 200\text{ oranges}\\  \textbf{(D)}\ \text{an infinite number of oranges}\qquad \textbf{(E)}\ \text{none of these}$


	\item (\textit{AHSME 1953}) A refrigerator is offered at sale at $250.00 less successive discounts of 20% and 15%. The sale price of the refrigerator is:


 $\textbf{(A) } \text{35\% less than 250.00} \qquad \textbf{(B) } \text{65\% of 250.00} \qquad \textbf{(C) } \text{77\% of 250.00} \qquad  \textbf{(D) } \text{68\% of 250.00} \qquad \textbf{(E) } \text{none of these}$


	\item (\textit{AHSME 1953}) The factors of the expression $x^2+y^2$ are: 


 $\textbf{(A)}\ (x+y)(x-y) \qquad \textbf{(B)}\ (x+y)^2 \qquad \textbf{(C)}\ (x^{\frac{2}{3}}+y^{\frac{2}{3}})(x^{\frac{4}{3}}+y^{\frac{4}{3}})\\  \textbf{(D)}\ (x+iy)(x-iy)\qquad \textbf{(E)}\ \text{none of these}$


	\item (\textit{AHSME 1953}) The roots of $x(x^2+8x+16)(4-x)=0$ are: 


 $\textbf{(A)}\ 0 \qquad \textbf{(B)}\ 0,4 \qquad \textbf{(C)}\ 0,4,-4 \qquad \textbf{(D)}\ 0,4,-4,-4 \qquad \textbf{(E)}\ \text{none of these}$


	\item (\textit{AHSME 1953}) If $\log_6 x=2.5$, the value of $x$ is: 


 $\textbf{(A)}\ 90 \qquad \textbf{(B)}\ 36 \qquad \textbf{(C)}\ 36\sqrt{6} \qquad \textbf{(D)}\ 0.5 \qquad \textbf{(E)}\ \text{none of these}$


	\item (\textit{AHSME 1953}) Charles has $5q + 1$ quarters and Richard has $q + 5$ quarters. The difference in their money in dimes is: 


 $\textbf{(A)}\ 10(q - 1) \qquad \textbf{(B)}\ \frac {2}{5}(4q - 4) \qquad \textbf{(C)}\ \frac {2}{5}(q - 1) \\  \textbf{(D)}\ \frac{5}{2}(q-1)\qquad \textbf{(E)}\ \text{none of these}$


	\item (\textit{AHSME 1953}) The fraction $\frac{\sqrt{a^2+x^2}-\frac{x^2-a^2}{\sqrt{a^2+x^2}}}{a^2+x^2}$ reduces to: 


 $\textbf{(A)}\ 0 \qquad \textbf{(B)}\ \frac{2a^2}{a^2+x^2} \qquad \textbf{(C)}\ \frac{2x^2}{(a^2+x^2)^{\frac{3}{2}}}\qquad \textbf{(D)}\ \frac{2a^2}{(a^2+x^2)^{\frac{3}{2}}}\qquad \textbf{(E)}\ \frac{2x^2}{a^2+x^2}$


	\item (\textit{AHSME 1953}) The value of $x$ at the intersection of $y=\frac{8}{x^2+4}$ and $x+y=2$ is: 


 $\textbf{(A)}\ -2+\sqrt{5} \qquad \textbf{(B)}\ -2-\sqrt{5} \qquad \textbf{(C)}\ 0 \qquad \textbf{(D)}\ 2 \qquad \textbf{(E)}\ \text{none of these}$


	\item (\textit{AHSME 1953}) The number of ounces of water needed to reduce $9$ ounces of shaving lotion containing $50$ \% alcohol to a lotion containing $30$ \% alcohol is: 


 $\textbf{(A)}\ 3 \qquad \textbf{(B)}\ 4 \qquad \textbf{(C)}\ 5 \qquad \textbf{(D)}\ 6 \qquad \textbf{(E)}\ 7$


	\item (\textit{AHSME 1953}) The number of revolutions of a wheel, with fixed center and with an outside diameter of $6$ feet, required to cause a point on the rim to go one mile is: 


 $\textbf{(A)}\ 880 \qquad \textbf{(B)}\ \frac{440}{\pi} \qquad \textbf{(C)}\ \frac{880}{\pi} \qquad \textbf{(D)}\ 440\pi\qquad \textbf{(E)}\ \text{none of these}$


	\item (\textit{AHSME 1953}) A running track is the ring formed by two concentric circles. It is $10$ feet wide. The circumference of the two circles differ by about: 


 $\textbf{(A)}\ 10\text{ feet} \qquad \textbf{(B)}\ 30\text{ feet} \qquad \textbf{(C)}\ 60\text{ feet} \qquad \textbf{(D)}\ 100\text{ feet}\\ \textbf{(E)}\ \text{none of these}$


	\item (\textit{AHSME 1953}) The diameters of two circles are $8$ inches and $12$ inches respectively. The ratio of the area of the smaller to the area of the larger circle is: 


 $\textbf{(A)}\ \frac{2}{3} \qquad \textbf{(B)}\ \frac{4}{9} \qquad \textbf{(C)}\ \frac{9}{4} \qquad \textbf{(D)}\ \frac{1}{2}\qquad \textbf{(E)}\ \text{none of these}$


	\item (\textit{AHSME 1953}) A triangle and a trapezoid are equal in area. They also have the same altitude. If the base of the triangle is 18 inches, the median of the trapezoid is: 


 $\textbf{(A)}\ 36\text{ inches} \qquad \textbf{(B)}\ 9\text{ inches} \qquad \textbf{(C)}\ 18\text{ inches}\\  \textbf{(D)}\ \text{not obtainable from these data}\qquad \textbf{(E)}\ \text{none of these}$


	\item (\textit{AHSME 1953}) Given the larger of two circles with center $P$ and radius $p$ and the smaller with center $Q$ and radius $q$. Draw $PQ$. Which of the following statements is false? 


 $\textbf{(A)}\ p-q\text{ can be equal to }\overline{PQ}\\  \textbf{(B)}\ p+q\text{ can be equal to }\overline{PQ}\\  \textbf{(C)}\ p+q\text{ can be less than }\overline{PQ}\\  \textbf{(D)}\ p-q\text{ can be less than }\overline{PQ}\\ \textbf{(E)}\ \text{none of these}$


	\item (\textit{AHSME 1953}) A circular piece of metal of maximum size is cut out of a square piece and then a square piece of maximum size is cut out of the circular piece. The total amount of metal wasted is: 


 $\textbf{(A)}\ \frac{1}{4} \text{ the area of the original square}\\  \textbf{(B)}\ \frac{1}{2}\text{ the area of the original square}\\ \textbf{(C)}\ \frac{1}{2}\text{ the area of the circular piece}\\ \textbf{(D)}\ \frac{1}{4}\text{ the area of the circular piece}\\ \textbf{(E)}\ \text{none of these}$


	\item (\textit{AHSME 1953}) Adams plans a profit of $10$ \% on the selling price of an article and his expenses are $15$ \% of sales. The rate of markup on an article that sells for $ $5.00$ is: 


 $\textbf{(A)}\ 20\% \qquad \textbf{(B)}\ 25\% \qquad \textbf{(C)}\ 30\% \qquad \textbf{(D)}\ 33\frac {1}{3}\% \qquad \textbf{(E)}\ 35\%$


	\item (\textit{AHSME 1953}) A man has part of $ $4500$ invested at $4$ % and the rest at $6$ %. If his annual return on each investment is the same, the average rate of interest which he realizes of the $4500 is: 


 $\textbf{(A)}\ 5\% \qquad \textbf{(B)}\ 4.8\% \qquad \textbf{(C)}\ 5.2\% \qquad \textbf{(D)}\ 4.6\% \qquad \textbf{(E)}\ \text{none of these}$


	\item (\textit{AHSME 1953}) One of the factors of $x^4+4$ is: 


 $\textbf{(A)}\ x^2+2 \qquad \textbf{(B)}\ x+1 \qquad \textbf{(C)}\ x^2-2x+2 \qquad \textbf{(D)}\ x^2-4\\  \textbf{(E)}\ \text{none of these}$


	\item (\textit{AHSME 1953}) In the expression $xy^2$, the values of $x$ and $y$ are each decreased $25$ \%; the value of the expression is: 


 $\textbf{(A)}\ \text{decreased } 50\% \qquad \textbf{(B)}\ \text{decreased }75\%\\  \textbf{(C)}\ \text{decreased }\frac{37}{64}\text{ of its value}\qquad \textbf{(D)}\ \text{decreased }\frac{27}{64}\text{ of its value}\\ \textbf{(E)}\ \text{none of these}$


	\item (\textit{AHSME 1953}) If $y=x+\frac{1}{x}$, then $x^4+x^3-4x^2+x+1=0$ becomes: 


 $\textbf{(A)}\ x^2(y^2+y-2)=0 \qquad \textbf{(B)}\ x^2(y^2+y-3)=0\\  \textbf{(C)}\ x^2(y^2+y-4)=0 \qquad \textbf{(D)}\ x^2(y^2+y-6)=0\\ \textbf{(E)}\ \text{none of these}$


	\item (\textit{AHSME 1953}) If $\log_{10} (x^2-3x+6)=1$, the value of $x$ is: 


 $\textbf{(A)}\ 10\text{ or }2 \qquad \textbf{(B)}\ 4\text{ or }-2 \qquad \textbf{(C)}\ 3\text{ or }-1 \qquad \textbf{(D)}\ 4\text{ or }-1\\ \textbf{(E)}\ \text{none of these}$


	\item (\textit{AHSME 1953}) The logarithm of $27\sqrt[4]{9}\sqrt[3]{9}$ to the base $3$ is: 


 $\textbf{(A)}\ 8\frac{1}{2} \qquad \textbf{(B)}\ 4\frac{1}{6} \qquad \textbf{(C)}\ 5 \qquad \textbf{(D)}\ 3 \qquad \textbf{(E)}\ \text{none of these}$


	\item (\textit{AHSME 1953}) The equation $\sqrt {x + 10} - \frac {6}{\sqrt {x + 10}} = 5$ has: 


 $\textbf{(A)}\ \text{an extraneous root between } - 5\text{ and } - 1 \\  \textbf{(B)}\ \text{an extraneous root between }-10\text{ and }-6\\ \textbf{(C)}\ \text{a true root between }20\text{ and }25\qquad \textbf{(D)}\ \text{two true roots}\\ \textbf{(E)}\ \text{two extraneous roots}$


	\item (\textit{AHSME 1953}) If $a,b,c$ are positive integers less than $10$, then $(10a + b)(10a + c) = 100a(a + 1) + bc$ if: 


 $\textbf{(A)}\ b + c = 10 \qquad \textbf{(B)}\ b = c \qquad \textbf{(C)}\ a + b = 10 \qquad \textbf{(D)}\ a = b \\  \textbf{(E)}\ a+b+c = 10$


	\item (\textit{AHSME 1953}) In a geometric progression whose terms are positive, any term is equal to the sum of the next two following terms. then the common ratio is: 


 $\textbf{(A)}\ 1 \qquad \textbf{(B)}\ \text{about }\frac{\sqrt{5}}{2} \qquad \textbf{(C)}\ \frac{\sqrt{5}-1}{2}\qquad \textbf{(D)}\ \frac{1-\sqrt{5}}{2}\qquad \textbf{(E)}\ \frac{2}{\sqrt{5}}$


	\item (\textit{AHSME 1953}) The base of a triangle is $15$ inches. Two lines are drawn parallel to the base, terminating in the other two sides, and dividing the triangle into three equal areas. The length of the parallel closer to the base is: 


 $\textbf{(A)}\ 5\sqrt{6}\text{ inches} \qquad \textbf{(B)}\ 10\text{ inches} \qquad \textbf{(C)}\ 4\sqrt{3}\text{ inches}\qquad \textbf{(D)}\ 7.5\text{ inches}\\  \textbf{(E)}\ \text{none of these}$


	\item (\textit{AHSME 1953}) The radius of the first circle is $1$ inch, that of the second $\frac{1}{2}$ inch, that of the third $\frac{1}{4}$ inch and so on indefinitely. 
The sum of the areas of the circles is: 


 $\textbf{(A)}\ \frac{3\pi}{4} \qquad \textbf{(B)}\ 1.3\pi \qquad \textbf{(C)}\ 2\pi \qquad \textbf{(D)}\ \frac{4\pi}{3}\qquad \textbf{(E)}\ \text{none of these}$


	\item (\textit{AHSME 1953}) In $\triangle ABC$, sides $a,b$ and $c$ are opposite $\angle{A},\angle{B}$ and $\angle{C}$ respectively. $AD$ bisects $\angle{A}$ and meets $BC$ at $D$. 
Then if $x = \overline{CD}$ and $y = \overline{BD}$ the correct proportion is: 


 $\textbf{(A)}\ \frac {x}{a} = \frac {a}{b + c} \qquad \textbf{(B)}\ \frac {x}{b} = \frac {a}{a + c} \qquad \textbf{(C)}\ \frac{y}{c}=\frac{c}{b+c}\\ \textbf{(D)}\ \frac{y}{c}=\frac{a}{b+c}\qquad \textbf{(E)}\ \frac{x}{y}=\frac{c}{b}$


	\item (\textit{AHSME 1953}) The number of significant digits in the measurement of the side of a square whose computed area is $1.1025$ square inches to 
the nearest ten-thousandth of a square inch is: 


 $\textbf{(A)}\ 2 \qquad \textbf{(B)}\ 3 \qquad \textbf{(C)}\ 4 \qquad \textbf{(D)}\ 5 \qquad \textbf{(E)}\ 1$


	\item (\textit{AHSME 1953}) A house worth $ $9000$ is sold by Mr. A to Mr. B at a $10$ % loss. Mr. B sells the house back to Mr. A at a $10$ % gain. 
The result of the two transactions is: 


 $\textbf{(A)}\ \text{Mr. A breaks even} \qquad \textbf{(B)}\ \text{Mr. B gains }\$900 \qquad \textbf{(C)}\ \text{Mr. A loses }\$900\\ \textbf{(D)}\ \text{Mr. A loses }\$810\qquad \textbf{(E)}\ \text{Mr. B gains }\$1710$


	\item (\textit{AHSME 1953}) The rails on a railroad are $30$ feet long. As the train passes over the point where the rails are joined, there is an audible click. 
The speed of the train in miles per hour is approximately the number of clicks heard in: 


 $\textbf{(A)}\ 20\text{ seconds} \qquad \textbf{(B)}\ 2\text{ minutes} \qquad \textbf{(C)}\ 1\frac{1}{2}\text{ minutes}\qquad \textbf{(D)}\ 5\text{ minutes}\\ \textbf{(E)}\ \text{none of these}$


	\item (\textit{AHSME 1953}) Each angle of a rectangle is trisected. The intersections of the pairs of trisectors adjacent to the same side always form: 


 $\textbf{(A)}\ \text{a square} \qquad \textbf{(B)}\ \text{a rectangle} \qquad \textbf{(C)}\ \text{a parallelogram with unequal sides}\\ \textbf{(D)}\ \text{a rhombus}\qquad \textbf{(E)}\ \text{a quadrilateral with no special properties}$


	\item (\textit{AHSME 1953}) The perimeter of an isosceles right triangle is $2p$. Its area is: 


 $\textbf{(A)}\ (2+\sqrt{2})p \qquad \textbf{(B)}\ (2-\sqrt{2})p \qquad \textbf{(C)}\ (3-2\sqrt{2})p^2\\  \textbf{(D)}\ (1-2\sqrt{2})p^2\qquad \textbf{(E)}\ (3+2\sqrt{2})p^2$


	\item (\textit{AHSME 1953}) If one side of a triangle is $12$ inches and the opposite angle is $30^{\circ}$, then the diameter of the circumscribed circle is: 


 $\textbf{(A)}\ 18\text{ inches} \qquad \textbf{(B)}\ 30\text{ inches} \qquad \textbf{(C)}\ 24\text{ inches} \qquad \textbf{(D)}\ 20\text{ inches}\\ \textbf{(E)}\ \text{none of these}$


	\item (\textit{AHSME 1953}) If $f(x)=\frac{x(x-1)}{2}$, then $f(x+2)$ equals: 


 $\textbf{(A)}\ f(x)+f(2) \qquad \textbf{(B)}\ (x+2)f(x) \qquad \textbf{(C)}\ x(x+2)f(x) \qquad \textbf{(D)}\ \frac{xf(x)}{x+2}\\ \textbf{(E)}\ \frac{(x+2)f(x+1)}{x}$


	\item (\textit{AHSME 1953}) Determine $m$ so that $4x^2-6x+m$ is divisible by $x-3$. The obtained value, $m$, is an exact divisor of: 


 $\textbf{(A)}\ 12 \qquad \textbf{(B)}\ 20 \qquad \textbf{(C)}\ 36 \qquad \textbf{(D)}\ 48 \qquad \textbf{(E)}\ 64$


	\item (\textit{AHSME 1953}) The base of an isosceles triangle is $6$ inches and one of the equal sides is $12$ inches. 
The radius of the circle through the vertices of the triangle is: 


 $\textbf{(A)}\ \frac{7\sqrt{15}}{5} \qquad \textbf{(B)}\ 4\sqrt{3} \qquad \textbf{(C)}\ 3\sqrt{5} \qquad \textbf{(D)}\ 6\sqrt{3}\qquad \textbf{(E)}\ \text{none of these}$


	\item (\textit{AHSME 1953}) If $f(a)=a-2$ and $F(a,b)=b^2+a$, then $F(3,f(4))$ is: 


 $\textbf{(A)}\ a^2-4a+7 \qquad \textbf{(B)}\ 28 \qquad \textbf{(C)}\ 7 \qquad \textbf{(D)}\ 8 \qquad \textbf{(E)}\ 11$


	\item (\textit{AHSME 1953}) The product, $\log_a b \cdot \log_b a$ is equal to: 


 $\textbf{(A)}\ 1 \qquad \textbf{(B)}\ a \qquad \textbf{(C)}\ b \qquad \textbf{(D)}\ ab \qquad \textbf{(E)}\ \text{none of these}$


	\item (\textit{AHSME 1953}) The negation of the statement "all men are honest," is: 


 $\textbf{(A)}\ \text{no men are honest} \qquad \textbf{(B)}\ \text{all men are dishonest} \\  \textbf{(C)}\ \text{some men are dishonest}\qquad \textbf{(D)}\ \text{no men are dishonest}\\ \textbf{(E)}\ \text{some men are honest}$


	\item (\textit{AHSME 1953}) A girls' camp is located $300$ rods from a straight road. On this road, a boys' camp is located $500$ rods from the girls' camp. 
It is desired to build a canteen on the road which shall be exactly the same distance from each camp. 
The distance of the canteen from each of the camps is: 


 $\textbf{(A)}\ 400\text{ rods} \qquad \textbf{(B)}\ 250\text{ rods} \qquad \textbf{(C)}\ 87.5\text{ rods} \qquad \textbf{(D)}\ 200\text{ rods}\\ \textbf{(E)}\ \text{none of these}$


	\item (\textit{AHSME 1953}) The centers of two circles are $41$ inches apart. The smaller circle has a radius of $4$ inches and the larger one has a radius of $5$ inches. 
The length of the common internal tangent is: 


 $\textbf{(A)}\ 41\text{ inches} \qquad \textbf{(B)}\ 39\text{ inches} \qquad \textbf{(C)}\ 39.8\text{ inches} \qquad \textbf{(D)}\ 40.1\text{ inches}\\ \textbf{(E)}\ 40\text{ inches}$


	\item (\textit{AHSME 1953}) If the price of an article is increased by percent $p$, then the decrease in percent of sales must not exceed $d$ in order to yield the same income. 
The value of $d$ is: 


 $\textbf{(A)}\ \frac{1}{1+p} \qquad \textbf{(B)}\ \frac{1}{1-p} \qquad \textbf{(C)}\ \frac{p}{1+p} \qquad \textbf{(D)}\ \frac{p}{p-1}\qquad \textbf{(E)}\ \frac{1-p}{1+p}$


	\item (\textit{AHSME 1953}) In solving a problem that reduces to a quadratic equation one student makes a mistake only in the constant term of the equation and 
obtains $8$ and $2$ for the roots. Another student makes a mistake only in the coefficient of the first degree term and 
find $-9$ and $-1$ for the roots. The correct equation was: 


 $\textbf{(A)}\ x^2-10x+9=0 \qquad \textbf{(B)}\ x^2+10x+9=0 \qquad \textbf{(C)}\ x^2-10x+16=0\\  \textbf{(D)}\ x^2-8x-9=0\qquad \textbf{(E)}\ \text{none of these}$


	\item (\textit{AHSME 1953}) The lengths of two line segments are $a$ units and $b$ units respectively. Then the correct relation between them is: 


 $\textbf{(A)}\ \frac{a+b}{2} > \sqrt{ab} \qquad \textbf{(B)}\ \frac{a+b}{2} < \sqrt{ab} \qquad \textbf{(C)}\ \frac{a+b}{2}=\sqrt{ab}\\ \textbf{(D)}\ \frac{a+b}{2}\leq\sqrt{ab}\qquad \textbf{(E)}\ \frac{a+b}{2}\geq\sqrt{ab}$


	\item (\textit{AHSME 1953}) Instead of walking along two adjacent sides of a rectangular field, a boy took a shortcut along the diagonal of the field and 
saved a distance equal to $\frac{1}{2}$ the longer side. The ratio of the shorter side of the rectangle to the longer side was: 


 $\textbf{(A)}\ \frac{1}{2} \qquad \textbf{(B)}\ \frac{2}{3} \qquad \textbf{(C)}\ \frac{1}{4} \qquad \textbf{(D)}\ \frac{3}{4}\qquad \textbf{(E)}\ \frac{2}{5}$


	\item (\textit{AHSME 1953}) If $x>0$, then the correct relationship is: 


 $\textbf{(A)}\ \log (1+x) = \frac{x}{1+x} \qquad \textbf{(B)}\ \log (1+x) < \frac{x}{1+x} \\  \textbf{(C)}\ \log(1+x) > x\qquad \textbf{(D)}\ \log (1+x) < x\qquad \textbf{(E)}\ \text{none of these}$


	\item (\textit{AHSME 1953}) If the larger base of an isosceles trapezoid equals a diagonal and the smaller base equals the altitude, 
then the ratio of the smaller base to the larger base is: 


 $\textbf{(A)}\ \frac{1}{2} \qquad \textbf{(B)}\ \frac{2}{3} \qquad \textbf{(C)}\ \frac{3}{4} \qquad \textbf{(D)}\ \frac{3}{5}\qquad \textbf{(E)}\ \frac{2}{5}$


	\item (\textit{AHSME 1953}) The coordinates of $A,B$ and $C$ are $(5,5),(2,1)$ and $(0,k)$ respectively. 
The value of $k$ that makes $\overline{AC}+\overline{BC}$ as small as possible is: 


 $\textbf{(A)}\ 3 \qquad \textbf{(B)}\ 4\frac{1}{2} \qquad \textbf{(C)}\ 3\frac{6}{7} \qquad \textbf{(D)}\ 4\frac{5}{6}\qquad \textbf{(E)}\ 2\frac{1}{7}$


	\item (\textit{AHSME 1953}) One of the sides of a triangle is divided into segments of $6$ and $8$ units by the point of tangency of the inscribed circle. 
If the radius of the circle is $4$, then the length of the shortest side of the triangle is: 


 $\textbf{(A)}\ 12\text{ units} \qquad \textbf{(B)}\ 13\text{ units} \qquad \textbf{(C)}\ 14\text{ units} \qquad \textbf{(D)}\ 15\text{ units}\qquad \textbf{(E)}\ 16\text{ units}$


\end{enumerate}